\section{A disjoint sparse partition of bounded-width CNFs}
\label{sec:sparse}

We next reduce arbitrary bounded-width CNF tests to tests with only linearly
many clauses.  The proof adapts the sparsification method of Impagliazzo,
Paturi, and Zane~\cite[Section~2]{IPZ01}, including its occurrence bounds,
insertion accounting, and entropy count of branches.  A disjoint
refinement is described by Dell, Husfeldt, Marx, Taslaman, and
Wahl\'en~\cite[Appendix~A of the full version]{DHMTW14}.
Here we prove a version with the bookkeeping needed for our partition.
Disjointness
lets us add the correlations of the pieces without inclusion--exclusion.

\begin{lemma}[Sparse partition]\label{lem:sparse-partition}
For every integer $b\ge 1$ and real $\eta>0$, there is a positive integer
$M=M(b,\eta)$ with the following property.  Every CNF indicator $H$ of width
at most $b$ on $m\ge 1$ variables can be written pointwise as
\[
  H=\sum_{a=1}^{q} H_a,
  \qquad q\le 2^{\eta m},
\]
where the $H_a$ are CNF indicators of width at most $b$, each represented by
at most $Mm$ clauses, and their satisfying sets are pairwise disjoint.
The empty sum is permitted when $H=0$.
\end{lemma}

We use the classical sunflower bound of Erd\H{o}s and Rado~\cite{ER60}
to identify a branching step, and include its short proof.
A \emph{sunflower} is a collection of distinct sets whose pairwise
intersections all equal a fixed set $C$, called its \emph{core}.  The sets
obtained by removing $C$ are its \emph{petals}; they are pairwise disjoint.

\begin{lemma}\label{lem:sunflower-bound}
Let $s\ge 0$ and $r\ge 2$ be integers.  A family of more than
$s!(r-1)^s$ distinct sets of cardinality $s$ contains a sunflower with
$r$ members.
\end{lemma}

\begin{proof}
For $s=0$ there is at most one distinct set, so the hypothesis is impossible.
Proceed by induction on $s$.  If the family contains $r$ pairwise disjoint
sets, these already form a sunflower with empty core.  Otherwise, take a
maximal pairwise disjoint subfamily.  It has at most $r-1$ members, and its
union $U$ has at most $s(r-1)$ elements.  Maximality implies that every set
in the original family meets $U$.  Some element of $U$ therefore belongs
to more than
\[
  \frac{s!(r-1)^s}{s(r-1)}=(s-1)!(r-1)^{s-1}
\]
members.  Delete this element from all those members.  The resulting
sets remain distinct and have cardinality $s-1$.  By induction, $r$ of
them form a sunflower.  Restoring the deleted element gives the required
sunflower in the original family.
\end{proof}

\begin{proof}[Proof of Lemma~\ref{lem:sparse-partition}]
We regard clauses as sets of literals.  Discard true and tautological
clauses, duplicates, and any clause that contains another clause.  If an
empty clause remains, use the empty partition and stop.  Otherwise the
remaining clauses are nonempty and form an antichain under inclusion.
The empty family is allowed and represents the constant true function.

The partition will come from a binary branching tree.  Fix for now
arbitrary integers $r_2,\ldots,r_b\ge2$; we will choose their values after
bounding the number of leaves.  A sunflower is \emph{eligible at size $i$}
if it consists of $r_i$ size-$i$ clauses and has nonempty core.  All set
operations below concern literals.  In particular, disjoint petals may
contain opposite literals of one variable, and a branch may be
unsatisfiable.

\medskip\noindent
\textbf{Branching and disjointness.}
Each state consists of an active antichain and a separate conjunction of
unit clauses, called its \emph{side condition}.  Initially the side
condition is true.  It is never used to select sunflowers or to delete
active clauses.  Whenever an eligible sunflower exists, choose one of
the smallest possible size $i$, with fixed deterministic tie-breaking,
and write its clauses as
\[
 C\cup P_1,\ldots,C\cup P_{r_i}.
\]
The core $C$ is nonempty.  The petals are nonempty, pairwise disjoint, and
all of the same size; their nonemptiness follows from distinctness and
equal size of the sunflower members.  Make two updates:
\begin{itemize}
 \item On branch $0$, insert the clause $C$ into the active family.
 \item On branch $1$, insert all petal clauses $P_1,\ldots,P_{r_i}$ and
 append to the side condition the requirement that every literal of $C$
 be false.
\end{itemize}
After either insertion, delete all subsumed active clauses to restore
the antichain.  Every inserted clause has size strictly less than $i$,
so width never increases; the side condition uses only unit clauses.
The extra requirement on branch $1$ is the disjointness device
of~\cite[Appendix~A of the full version]{DHMTW14}.  Keeping it separate
ensures that every deletion from the active family is caused by an
active insertion, a fact needed in the counting argument below.

Every clause in an insertion batch is new and remains active immediately
after the update.  Indeed, an inserted clause $D$ is a strict subset of
some current clause $A$.  A current clause $B\subseteq D$ would satisfy
$B\subsetneq A$, contrary to the antichain property.  Within a petal
batch, distinct equal-size clauses do not subsume one another either.

Let $F$ be the active CNF indicator before a step, and let $F_0,F_1$ be
the active indicators after its two updates, without side conditions.
Subsumption deletion preserves the conjunction, so
\[
 F_0=F\,\one[C\text{ is true}],\qquad
 F_1=F\prod_{j=1}^{r_i}\one[P_j\text{ is true}].
\]
When $C$ is false, the selected clauses in $F$ require every petal to be
true.  Therefore
\begin{equation}\label{eq:sparse-branch-identity}
 F=F_0+\one[C\text{ is false}]F_1,
\end{equation}
with disjoint supports on the right.  Multiplying by the accumulated
side condition gives the corresponding disjoint identity for the state.
Once we prove termination, iteration of this identity will give a
partition into leaf indicators.

\medskip\noindent
\textbf{Occurrence bounds and insertion budgets.}
We next bound the number of steps on every finite path prefix.  The
smallest-size rule bounds how many active clauses inserted in earlier
branching steps can contain any one literal.  Define
\[
 e_1=1,\qquad e_i=(i-1)!(r_i-1)^{i-1}\quad(2\le i\le b).
\]
If the current size-$i$ clauses contain no eligible sunflower, each
literal occurs in at most $e_i$ of them.  For $i\ge2$, remove a fixed
literal from the clauses containing it.  More than $e_i$ distinct
resulting sets would contain an $r_i$-member sunflower by
Lemma~\ref{lem:sunflower-bound}; restoring the literal gives an eligible
sunflower with nonempty core.  For $i=1$, the claim follows because the
clauses are distinct.

Call an active clause \emph{added} if it was inserted in a branching step.
At every state, for every size $u$ and every literal, we claim
\begin{equation}\label{eq:sparse-added-invariant}
 \#\{\text{active added size-$u$ clauses containing this literal}\}
 \le e_u+1.
\end{equation}
Initially there are no added clauses.  A step can increase this count
only by inserting size-$u$ clauses, and such a step processes a size
$i>u$.  The smallest-size rule then guarantees that there is no eligible
size-$u$ sunflower before the insertion.  Thus the occurrence bound
applies to all current size-$u$ clauses, giving at most $e_u$ occurrences
of each literal.  The insertion adds at most one further occurrence,
because it adds either one core or pairwise disjoint petals.  Deletions
only decrease the count.  This proves the invariant; the initial family
itself need not satisfy an occurrence bound.

For any finite path prefix, let $N_u$ count the size-$u$ insertions,
counting members of a batch individually.  Every deleted added size-$u$
clause was subsumed by a newly inserted strictly smaller clause.
Equality is excluded by the freshness of insertions, and side
conditions never delete active clauses.  Charge each such deletion to
one smaller inserted clause that caused it.  A fixed inserted clause
$D$ receives at most $e_u+1$ charges: choose a literal of the nonempty
clause $D$, and apply \eqref{eq:sparse-added-invariant} just before its
batch was inserted.  Every size-$u$ clause subsumed by $D$ contains that
literal.  If several batch members could receive a charge, choose only
one of them.

At the end of the prefix, at most $2m(e_u+1)$ added size-$u$ clauses are
still active, since each is nonempty and there are only $2m$ literals.
Counting deleted and surviving insertions therefore gives
\begin{equation}\label{eq:sparse-insertion-charge}
 N_u\le(e_u+1)\left(2m+\sum_{v<u}N_v\right).
\end{equation}
Define the positive integers
\begin{equation}\label{eq:sparse-constants}
 D_1=4,\qquad
 D_i=(e_i+1)\left(2+\sum_{u<i}D_u\right)\quad(2\le i\le b).
\end{equation}
The same recurrence holds for $i=1$ with an empty sum.  Induction on
$u$ in \eqref{eq:sparse-insertion-charge} proves
\begin{equation}\label{eq:sparse-insertion-budget}
 N_u\le D_um.
\end{equation}
These bounds hold on every finite prefix, without assuming termination.
Every step inserts at least one clause of size less than $b$, so every
path has at most $m\sum_{u<b}D_u$ steps.  The binary tree is therefore
finite.

For $2\le i\le b$, put $F_i=\sum_{u<i}D_u$.  Each step processing size
$i$ contributes at least one insertion below size $i$, and each
branch-$1$ step contributes exactly $r_i$ such insertions.  Hence
\begin{equation}\label{eq:sparse-branch-budgets}
 \#\{\text{steps processing size }i\}\le F_im,\qquad
 \#\{\text{branch-$1$ steps processing size }i\}\le F_im/r_i.
\end{equation}
Insertions caused by other processed sizes may also use this budget;
including them only enlarges the upper bound.

\medskip\noindent
\textbf{Counting leaves and choosing thresholds.}
For each leaf and each $i=2,\ldots,b$, record in order the decisions made
at steps processing size $i$, and pad this binary string with zeros to
length $F_im$.  By \eqref{eq:sparse-branch-budgets}, it has at most
$F_im/r_i$ ones.  The tuple of strings determines the leaf.  Indeed,
distinct leaf paths have a first differing branch: neither can be a
proper prefix of the other.  Before this difference their states agree,
so their selected size and their positions in that size's stream agree.
The different bits occur before padding in both streams, so the padded
tuples remain different.

Write $h_2(x)=-x\log_2x-(1-x)\log_2(1-x)$ for $0<x<1$.
For integer $N$ and $0<x\le1/2$, the number of binary strings of length
$N$ with at most $Nx$ ones is at most $2^{Nh_2(x)}$.  A string with
$j\le Nx$ ones has Bernoulli-$x$ probability
\[
 x^j(1-x)^{N-j}\ge x^{Nx}(1-x)^{N-Nx}=2^{-Nh_2(x)},
\]
because $x/(1-x)\le1$.  Their total probability is at most one, proving
the bound even when $Nx$ is not an integer.  Applied to the decision
streams, this gives
\[
 \#\{\text{leaves}\}\le
 \prod_{i=2}^{b}2^{mF_i h_2(1/r_i)}.
\]
We can now choose the thresholds to make this exponent small.  The
constant $F_i$ depends only on $r_2,\ldots,r_{i-1}$, because it involves
only $D_u$ with $u<i$.  Choose $r_2,r_3,\ldots,r_b$ in order, each large
enough that
\begin{equation}\label{eq:sparse-thresholds}
 F_i h_2(1/r_i)\le\eta/b,\qquad r_i\in\N,\quad r_i\ge2.
\end{equation}
This is possible since $h_2(1/r)\to0$.  An explicit choice is
$r_i=2^{q_i}$, where $q_i\ge1$ is the least integer satisfying
$F_i(q_i+2)2^{-q_i}\le\eta/b$: the inequality
$-(1-x)\ln(1-x)\le x$ implies
$h_2(2^{-q})\le(q+\log_2e)2^{-q}<(q+2)2^{-q}$.
For these choices,
\begin{equation}\label{eq:sparse-leaf-count}
 \#\{\text{leaves}\}\le
 \prod_{i=2}^{b}2^{mF_i h_2(1/r_i)}
 \le2^{\eta m(b-1)/b}\le2^{\eta m}.
\end{equation}

\medskip\noindent
\textbf{Clause counts at the leaves.}
At a terminal state no eligible sunflower remains.  The occurrence
bound now applies to the entire active antichain: each literal occurs
in at most $e_i$ size-$i$ clauses, so there are at most $2me_i$ such
clauses.  The side condition contributes at most $2m$ distinct unit
clauses after repetitions are removed.  Thus the leaf indicator has
width at most $b$ and at most $Mm$ clauses, where
\begin{equation}\label{eq:sparse-M}
 M=2+2\sum_{i=1}^{b}e_i.
\end{equation}
All constants depend only on $b,\eta$.  An inconsistent side condition
gives an empty piece, which may be discarded.  The branching identities
\eqref{eq:sparse-branch-identity} give coverage and disjointness, and
\eqref{eq:sparse-leaf-count} bounds the number of pieces.

When $b=1$ there are no branching sizes or streams.  The tree has one
leaf, its empty product of stream counts is one, and at most $2m$
distinct unit clauses suffice; \eqref{eq:sparse-M} gives $M=4$.
The constant true function has the empty conjunction representation,
and the constant false function can use the empty partition.  This
covers every $m\ge1$.
\end{proof}
